\chapter{Metodologia de comparação}
\label{cap:comparacao}

Uma comparação de controladores só é defensável se cada um for avaliado \textbf{no seu
melhor ajuste}, para o \textbf{mesmo objetivo}, nos \textbf{mesmos cenários}, e com
estatística sobre várias execuções. Este capítulo propõe um protocolo.

\section{Métricas}
O nó publica \code{PursuitMetrics} a cada ciclo (\code{control/performance\_metrics.cpp}).
Todas as integrais são em tempo, então ficam robustas ao \emph{jitter} do laço:
\begin{table}[H]
\centering
\begin{tabular}{@{}lll@{}}
\toprule
Métrica & Definição & Significado \\
\midrule
\code{rmse\_position} & $\sqrt{\frac1T\int\lVert\vect p_{des}\rVert^2dt}$ & precisão do standoff \\
\code{rmse\_velocity} & $\sqrt{\frac1T\int\lVert\vect v_c-\vect v_t\rVert^2dt}$ & casamento de velocidade \\
\code{rmse\_los} & $\sqrt{\frac1T\int\psi_{LOS}^2dt}$ & alvo centrado na câmera (yaw) \\
\code{max\_position\_error} & $\max\lVert\vect p_{des}\rVert$ & pior caso \\
\code{effort\_l2} / \code{effort\_l1} & $\int\lVert\vect a\rVert^2dt$ / $\int\lVert\vect a\rVert dt$ & energia / consumo \\
\code{solve\_time\_us} & tempo de cálculo da lei & custo computacional \\
--- & tempo até $\lVert\vect p_{des}\rVert<\SI{0,5}{m}$ & rapidez de engajamento \\
--- & fração do tempo com alvo rastreado & robustez (perda de alvo) \\
\bottomrule
\end{tabular}
\end{table}

\section{Cenários}
Use o \code{scenario\_runner} (pasta \code{scenarios/}) para que o alvo execute
exatamente a mesma trajetória em todos os ensaios. Proposta mínima:
\begin{enumerate}
  \item \textbf{S1 -- Regulação:} alvo pairando, perseguidor começando a 8--\SI{10}{m}
        (degrau de \SI{3}{m} no erro). Mede $t_s$, $M_p$ e esforço.
  \item \textbf{S2 -- Reta:} alvo a 1,0 e \SI{1,5}{m/s} constantes. Mede o erro de
        regime e o feedforward.
  \item \textbf{S3 -- Curva:} taxa de curva constante (\SI{0,3}{rad/s}). Mede o erro de
        modelo do alvo (velocidade constante).
  \item \textbf{S4 -- Para-e-anda:} acelerações bruscas do alvo. Mede rejeição de
        perturbação (onde SMC e LQI devem se destacar).
  \item \textbf{S5 -- Evasivo:} política \code{flee} do \code{training\_arena}.
  \item \textbf{S6 -- Saturação:} alvo mais rápido que \code{max\_vel}. Comportamento
        degradado.
\end{enumerate}
Cruze os cenários com duas condições de estimação: (a) \code{use\_tf\_tracking: true}
(\emph{ground truth}, isola o controlador) e (b) visão real (controlador + percepção).

\section{Ajuste justo: um custo comum}
Defina \emph{um} custo e otimize os parâmetros $\theta$ de \emph{cada} controlador sobre
um conjunto de cenários de \textbf{treino}:
\begin{equation}
  J(\theta)=\sum_{c\in\text{treino}}
  \left[\frac{1}{T_c}\int_0^{T_c}\!\left(\lVert\vect p_{des}\rVert^2
  + w_v\lVert\vect v_c-\vect v_t\rVert^2 + \rho\lVert\vect a\rVert^2\right)dt\right]
  + \mu\,(\ang{45}-\text{PM}(\theta))_+^2 .
  \label{eq:jcomum}
\end{equation}
Com $w_v=0{,}2$ e $\rho=0{,}05$ ($\approx$ as razões $q_v/q_p$ e $r/q_p$ do \code{yaml}),
o LQR com $Q=\mathrm{diag}(10,2)$ e $R=0{,}5$ é o ótimo teórico do modelo ideal: é a
\emph{linha de base}. Avalie depois nos cenários de \textbf{teste}, diferentes dos de
treino, para não favorecer quem ``decorou'' o cenário.

\begin{receita}
\textbf{Pipeline recomendado.}
\begin{enumerate}
  \item \textbf{MATLAB} (minutos): \code{otimizar\_ganhos.m} ajusta cada controlador em
        \code{sim\_perseguicao.m} (modelo PX4 + atraso + ruído + KF).
  \item \textbf{MATLAB} (Monte Carlo): \code{comparar\_controladores.m} roda 30
        sementes por cenário e gera tabela e boxplots.
  \item \textbf{Gazebo SITL}: copie os ganhos para o \code{yaml}, rode os cenários com o
        \code{session\_recorder} (um MCAP por perseguição) e ao menos 10 repetições.
  \item \textbf{Análise}: exporte \code{/pursuit/metrics} dos bags e compare com o
        MATLAB. A diferença entre simulação simplificada e SITL também é um resultado
        (fidelidade do modelo).
\end{enumerate}
\end{receita}

\section{Estatística}
\begin{itemize}
  \item Reporte média $\pm$ desvio-padrão e mediana [IQR] de cada métrica, por cenário.
  \item As distribuições raramente são normais: use \textbf{Kruskal--Wallis}
        (\code{kruskalwallis}) entre controladores e \textbf{Wilcoxon} pareado
        (\code{signrank}) entre pares, com a mesma semente para os dois.
  \item Gráficos: boxplot por métrica e uma ``fronteira de Pareto'' (RMSE de posição
        $\times$ esforço L2). Os pontos ótimos de LQR/MPC para vários $\rho$ formam a
        curva de referência.
\end{itemize}

\section{Modelo de tabela de resultados}
\begin{table}[H]
\centering
\caption{Exemplo de tabela (preencher com \code{comparar\_controladores.m}).}
\begin{tabular}{@{}lcccccc@{}}
\toprule
Controlador & RMSE pos [m] & RMSE vel [m/s] & máx. [m] & esforço L2 & $t_{0,5}$ [s] & solve [\si{\micro s}] \\
\midrule
PID PBVS & & & & & & \\
LQR & & & & & & \\
LQI & & & & & & \\
ASMC & & & & & & \\
MPC cinemático & & & & & & \\
MPC dinâmico & & & & & & \\
Hermite & & & & & & \\
PN/PP & & & & & & \\
\bottomrule
\end{tabular}
\end{table}
